% Part: first-order-logic
% Chapter: sequent-calculus
% Section: soundness.tex

\documentclass[../../../include/open-logic-section]{subfiles}

\begin{document}

\iftag{FOL}
      {\olfileid{fol}{seq}{sou}}
      {\olfileid{pl}{seq}{sou}}

\olsection{યથાર્થતા}

\begin{explain}
સિક્વન્ટ કલન જેવું !!^a{derivation} તંત્ર ખરેખર સધાતી ન હોય એવી વસ્તુઓ
!!{derive} ન કરી શકે, તો તે \emph{યથાર્થ} છે. તેથી યથાર્થતા એ
!!{derivation} તંત્રો માટે ખાતરી અપાવતો એક પ્રકારનો સુરક્ષા-ગુણધર્મ છે.
કયો સાબિતી-સૈદ્ધાંતિક ગુણધર્મ વિચારાધીન છે તેના આધારે, દાખલા તરીકે,
આપણે નીચેની બાબતો જાણવા માગીએ છીએ:
\begin{enumerate}
\item દરેક !!{derivable}~$!A$ એ \iftag{FOL}{પ્રામાણ્ય}{પુનરુક્તિ} છે;
\item જો !!a{sentence} અમુક બીજાં વાક્યોમાંથી !!{derivable} હોય, તો તે
  તેમનું ફલિત પણ છે;
\item જો !!{sentence}sનો ગણ અસુસંગત હોય, તો તે અસંતોષ્ય છે.
\end{enumerate}
!!a{derivation} તંત્રના આ મહત્વના ગુણધર્મો છે. તેમાંથી કોઈ ગુણધર્મ ન
સધાય, તો !!{derivation} તંત્ર ખામીયુક્ત છે—તે જરૂર કરતાં વધારે
!!{derive} કરશે. તેથી !!{derivation} તંત્રની યથાર્થતા સ્થાપવી સૌથી વધુ
મહત્વ ધરાવે છે.

આ બધા સાબિતી-સૈદ્ધાંતિક ગુણધર્મો સિક્વન્ટ કલનમાં અમુક સિક્વન્ટોની
!!{derivability} વડે વ્યાખ્યાયિત થાય છે. તેથી ઉપરનાં (1)--(3) સાબિત કરવા
!!{derivable} સિક્વન્ટોના અર્થાનુસારી ગુણધર્મો વિશે કંઈક સાબિત કરવું
જરૂરી છે. પહેલાં સિક્વન્ટ \emph{પ્રામાણ્ય} હોવાનો અર્થ વ્યાખ્યાયિત કરીશું
અને પછી દરેક !!{derivable} સિક્વન્ટ પ્રામાણ્ય છે તેમ બતાવીશું. ત્યાર બાદ
(1)--(3) આ પરિણામનાં અનુપ્રમેયો તરીકે મળે છે.
\end{explain}

\begin{defn}
\iftag{FOL}{!!^a{structure}~$\Struct M$}{!!^a{valuation}~$\pAssign{v}$}
સિક્વન્ટ $\Gamma \Sequent \Delta$ને \emph{સંતોષે} છે જો અને માત્ર જો
કોઈ $!A \in \Gamma$ માટે
$\iftag{FOL}{\Sat/{M}}{\pSat/{v}}{!A}$ હોય અથવા કોઈ
$!A \in \Delta$ માટે $\iftag{FOL}{\Sat{M}}{\pSat{v}}{!A}$ હોય.

દરેક \iftag{FOL}{!!{structure}~$\Struct M$}{!!{valuation}~$\pAssign{v}$}
સિક્વન્ટને સંતોષે જો અને માત્ર જો તે સિક્વન્ટ \emph{પ્રામાણ્ય} છે.
\end{defn}

\begin{thm}[યથાર્થતા]
\ollabel{thm:sequent-soundness} જો $\Log{LK}$ એ
$\Theta \Sequent \Xi$ને !!{derive}s કરે, તો $\Theta \Sequent \Xi$
પ્રામાણ્ય છે.
\end{thm}

\begin{proof}
$\Theta \Sequent \Xi$ની કોઈ !!a{derivation}~$\pi$ લો. $\pi$માંનાં
અનુમાનોની સંખ્યા~$n$ પર આગમનથી આગળ વધીએ.

અનુમાનોની સંખ્યા~$0$ હોય, તો $\pi$માં માત્ર આરંભિક સિક્વન્ટ છે. દરેક
આરંભિક સિક્વન્ટ $!A \Sequent !A$ સ્પષ્ટ રીતે પ્રામાણ્ય છે, કારણ કે દરેક
\iftag{FOL}{$\Struct M$}{$\pAssign{v}$} માટે કાં તો
$\iftag{FOL}{\Sat/{M}}{\pSat/{v}}{!A}$ અથવા
$\iftag{FOL}{\Sat{M}}{\pSat{v}}{!A}$.

અનુમાનોની સંખ્યા~0 કરતાં મોટી હોય, તો સૌથી નીચેના અનુમાનના પ્રકાર
પ્રમાણે કિસ્સા પાડીએ. આગમન-પરિકલ્પના મુજબ તે અનુમાનનાં આધાર-સિક્વન્ટો
પ્રામાણ્ય છે તેમ માની શકીએ, કારણ કે કોઈપણ આધાર-સિક્વન્ટની
!!{derivation}માંનાં અનુમાનોની સંખ્યા~$n$ કરતાં નાની છે.

પહેલાં માત્ર એક આધાર-સિક્વન્ટ ધરાવતાં સંભવિત અનુમાનો જોઈએ.
\begin{enumerate}
\item છેલ્લું અનુમાન શિથિલીકરણ છે. ત્યારે $\Theta \Sequent \Xi$ કાં તો
  $!A, \Gamma \Sequent \Delta$ છે (છેલ્લું અનુમાન
  \LeftR{\Weakening} હોય ત્યારે) અથવા
  $\Gamma \Sequent \Delta, !A$ છે (તે \RightR{\Weakening} હોય ત્યારે),
  અને !!{derivation}નો અંત નીચેનામાંથી કોઈ એક છે:
  \begin{prooftree}
    \AxiomC{}
    \Deduce$\Gamma \fCenter \Delta$
    \RightLabel{\LeftR{\Weakening}}
    \UnaryInf$!A, \Gamma \fCenter \Delta$
    \DisplayProof\qquad\bottomAlignProof
    \AxiomC{}
    \Deduce$\Gamma \fCenter \Delta$
    \RightLabel{\RightR{\Weakening}}
    \UnaryInf$\Gamma \fCenter \Delta, !A$
  \end{prooftree}
  આગમન-પરિકલ્પના મુજબ $\Gamma \Sequent \Delta$ પ્રામાણ્ય છે. એટલે કે,
  દરેક
  \iftag{FOL}{!!{structure}~$\Struct{M}$}{!!{valuation}~$\pAssign{v}$}
  માટે કાં તો કોઈ $!C \in \Gamma$ એવું છે કે
  $\iftag{FOL}{\Sat/{M}}{\pSat/{v}}{!C}$, અથવા કોઈ $!C \in \Delta$
  એવું છે કે $\iftag{FOL}{\Sat{M}}{\pSat{v}}{!C}$.

  જો કોઈ $!C \in \Gamma$ માટે
  $\iftag{FOL}{\Sat/{M}}{\pSat/{v}}{!C}$ હોય, તો ડાબા શિથિલીકરણમાં
  $\Theta = !A, \Gamma$ અને જમણા શિથિલીકરણમાં $\Theta = \Gamma$;
  તેથી બંને કિસ્સામાં કોઈ $!C \in \Theta$ માટે
  $\iftag{FOL}{\Sat/{M}}{\pSat/{v}}{!C}$ છે. એ જ રીતે, જો કોઈ
  $!C \in \Delta$ માટે $\iftag{FOL}{\Sat{M}}{\pSat{v}}{!C}$ હોય,
  તો ડાબા શિથિલીકરણમાં $\Xi = \Delta$ અને જમણા શિથિલીકરણમાં
  $\Xi = \Delta, !A$; તેથી બંને કિસ્સામાં કોઈ $!C \in \Xi$ માટે
  $\iftag{FOL}{\Sat{M}}{\pSat{v}}{!C}$ છે. પરિણામે
  $\Theta \Sequent \Xi$ પ્રામાણ્ય છે.
  \sourcecorrection{OLSEQ-004}{મૂળની \texttt{soundness.tex}માં બંને
  શિથિલીકરણના કિસ્સા એકસાથે લીધા પછી $\Theta = !A,\Gamma$ લખાયું છે,
  જે માત્ર ડાબા શિથિલીકરણ માટે સાચું છે. અહીં બંને કિસ્સામાં
  $\Theta$ અને $\Xi$ની ચોક્કસ કિંમતો લખી છે.}
\item છેલ્લું અનુમાન \LeftR{\lnot} છે. ત્યારે છેલ્લા અનુમાનનું
  આધાર-સિક્વન્ટ $\Gamma \Sequent \Delta, !A$ અને તેનું ફલિત-સિક્વન્ટ
  $\lnot !A, \Gamma \Sequent \Delta$ છે; એટલે કે !!{derivation}નો અંત
  આ પ્રમાણે છે:
  \begin{prooftree}
    \AxiomC{}
    \Deduce$\Gamma \fCenter \Delta, !A$
    \RightLabel{\LeftR{\lnot}}
    \UnaryInf$\lnot !A, \Gamma \fCenter \Delta$
  \end{prooftree}
  અને $\Theta = \lnot !A, \Gamma$, જ્યારે $\Xi = \Delta$.

  આગમન-પરિકલ્પના જણાવે છે કે $\Gamma \Sequent \Delta, !A$ પ્રામાણ્ય
  છે. એટલે કે, દરેક \iftag{FOL}{$\Struct{M}$}{$\pAssign{v}$} માટે કાં તો
  (a) કોઈ $!C \in \Gamma$ માટે
  $\iftag{FOL}{\Sat/{M}}{\pSat/{v}}{!C}$, અથવા (b) કોઈ
  $!C \in \Delta$ માટે $\iftag{FOL}{\Sat{M}}{\pSat{v}}{!C}$, અથવા
  (c) $\iftag{FOL}{\Sat{M}}{\pSat{v}}{!A}$. આપણે
  $\Theta \Sequent \Xi$ પણ પ્રામાણ્ય છે તેમ બતાવવા માગીએ છીએ. કોઈ
  \iftag{FOL}{!!a{structure}~$\Struct{M}$}{!!a{valuation}~$\pAssign{v}$}
  લો. (a) સધાય તો એવું $!C \in \Gamma$ છે કે
  $\iftag{FOL}{\Sat/{M}}{\pSat/{v}}{!C}$, અને $!C \in \Theta$ પણ છે.
  (b) સધાય તો એવું $!C \in \Delta$ છે કે
  $\iftag{FOL}{\Sat{M}}{\pSat{v}}{!C}$, અને $!C \in \Xi$ પણ છે. છેલ્લે,
  જો $\iftag{FOL}{\Sat{M}}{\pSat{v}}{!A}$ હોય, તો
  $\iftag{FOL}{\Sat/{M}}{\pSat/{v}}{\lnot !A}$. $\lnot !A \in \Theta$
  હોવાથી એવું $!C \in \Theta$ છે કે
  $\iftag{FOL}{\Sat/{M}}{\pSat/{v}}{!C}$. તેથી
  $\Theta \Sequent \Xi$ પ્રામાણ્ય છે.
\item છેલ્લું અનુમાન \RightR{\lnot} છે: અભ્યાસપ્રશ્ન.
\item છેલ્લું અનુમાન \LeftR{\land} છે. તેનાં બે રૂપ છે: ડાબી બાજુ
  $!A$ ધરાવતા આધાર-સિક્વન્ટમાંથી અથવા ડાબી બાજુ $!B$ ધરાવતા
  આધાર-સિક્વન્ટમાંથી $!A \land !B$નું અનુમાન થઈ શકે છે. પહેલા કિસ્સામાં
  $\pi$નો અંત આ પ્રમાણે છે:
  \begin{prooftree}
    \AxiomC{}
    \Deduce$!A, \Gamma \fCenter \Delta$
    \RightLabel{\LeftR{\land}}
    \UnaryInf$!A \land !B, \Gamma \fCenter \Delta$
  \end{prooftree}
  અને $\Theta = !A \land !B, \Gamma$, જ્યારે $\Xi = \Delta$. કોઈ
  \iftag{FOL}{!!a{structure}~$\Struct M$}{!!a{valuation}~$\pAssign{v}$}
  લો. આગમન-પરિકલ્પના મુજબ $!A, \Gamma \Sequent \Delta$ પ્રામાણ્ય
  હોવાથી કાં તો (a) $\iftag{FOL}{\Sat/{M}}{\pSat/{v}}{!A}$, અથવા
  (b) કોઈ $!C \in \Gamma$ માટે
  $\iftag{FOL}{\Sat/{M}}{\pSat/{v}}{!C}$, અથવા (c) કોઈ
  $!C \in \Delta$ માટે $\iftag{FOL}{\Sat{M}}{\pSat{v}}{!C}$.
  કિસ્સા (a)માં $\iftag{FOL}{\Sat/{M}}{\pSat/{v}}{!A \land !B}$;
  તેથી એવું $!C \in \Theta$ છે (એટલે કે $!A \land !B$) કે
  $\iftag{FOL}{\Sat/{M}}{\pSat/{v}}{!C}$. કિસ્સા (b)માં એવું
  $!C \in \Gamma$ છે કે $\iftag{FOL}{\Sat/{M}}{\pSat/{v}}{!C}$, અને
  $!C \in \Theta$ પણ છે. કિસ્સા (c)માં એવું $!C \in \Delta$ છે કે
  $\iftag{FOL}{\Sat{M}}{\pSat{v}}{!C}$, અને $\Xi = \Delta$ હોવાથી
  $!C \in \Xi$ પણ છે. આમ દરેક કિસ્સામાં
  \iftag{FOL}{$\Struct M$}{$\pAssign{v}$} એ
  $!A \land !B, \Gamma \Sequent \Delta$ને સંતોષે છે. આ
  \iftag{FOL}{$\Struct{M}$}{$\pAssign{v}$} મનસ્વી હતું, તેથી
  $!A \land !B, \Gamma \Sequent \Delta$ પ્રામાણ્ય છે. $!B$માંથી
  $!A \land !B$નું અનુમાન થતું હોય તે કિસ્સો $!A$ને $!B$થી બદલીને એ જ
  રીતે સંભાળાય છે.
  \sourcecorrection{OLSEQ-005}{મૂળની \texttt{soundness.tex}માં આ
  કિસ્સાના અંતે હમણાં સાબિત કરેલા ફલિત-સિક્વન્ટને બદલે
  $\Gamma \Sequent \Delta$ પ્રામાણ્ય હોવાનું લખાયું છે. અહીં નિયમનું
  સાચું ફલિત-સિક્વન્ટ $!A \land !B,\Gamma \Sequent \Delta$ પુનઃસ્થાપિત
  કર્યું છે.}
\item છેલ્લું અનુમાન \RightR{\lor} છે. તેનાં બે રૂપ છે: આધાર-સિક્વન્ટની
  જમણી બાજુના $!A$માંથી અથવા જમણી બાજુના $!B$માંથી $!A \lor !B$નું
  અનુમાન થઈ શકે છે. પહેલા કિસ્સામાં $\pi$નો અંત આ પ્રમાણે છે:
  \begin{prooftree}
    \AxiomC{}
    \Deduce$\Gamma \fCenter \Delta, !A$
    \RightLabel{\RightR{\lor}}
    \UnaryInf$\Gamma \fCenter \Delta, !A \lor !B$
  \end{prooftree}
  હવે $\Theta = \Gamma$ અને $\Xi = \Delta, !A \lor !B$. કોઈ
  \iftag{FOL}{!!a{structure}~$\Struct{M}$}{!!a{valuation}~$\pAssign{v}$}
  લો. $\Gamma \Sequent \Delta, !A$ પ્રામાણ્ય હોવાથી કાં તો
  (a) $\iftag{FOL}{\Sat{M}}{\pSat{v}}{!A}$, અથવા (b) કોઈ
  $!C \in \Gamma$ માટે $\iftag{FOL}{\Sat/{M}}{\pSat/{v}}{!C}$, અથવા
  (c) કોઈ $!C \in \Delta$ માટે $\iftag{FOL}{\Sat{M}}{\pSat{v}}{!C}$.
  કિસ્સા (a)માં $\iftag{FOL}{\Sat{M}}{\pSat{v}}{!A \lor !B}$.
  કિસ્સા (b)માં એવું $!C \in \Gamma$ છે કે
  $\iftag{FOL}{\Sat/{M}}{\pSat/{v}}{!C}$. કિસ્સા (c)માં એવું
  $!C \in \Delta$ છે કે $\iftag{FOL}{\Sat{M}}{\pSat{v}}{!C}$.
  આમ દરેક કિસ્સામાં \iftag{FOL}{$\Struct{M}$}{$\pAssign{v}$} એ
  $\Gamma \Sequent \Delta, !A \lor !B$, એટલે કે
  $\Theta \Sequent \Xi$, સંતોષે છે. આ
  \iftag{FOL}{$\Struct{M}$}{$\pAssign{v}$} મનસ્વી હતું, તેથી
  $\Theta \Sequent \Xi$ પ્રામાણ્ય છે. $!B$માંથી $!A \lor !B$નું
  અનુમાન થતું હોય તે કિસ્સો $!A$ને $!B$થી બદલીને એ જ રીતે સંભાળાય છે.
\item છેલ્લું અનુમાન \RightR{\lif} છે. ત્યારે $\pi$નો અંત આ પ્રમાણે છે:
  \begin{prooftree}
    \AxiomC{}
    \Deduce$!A, \Gamma \fCenter \Delta, !B$
    \RightLabel{\RightR{\lif}}
    \UnaryInf$\Gamma \fCenter \Delta, !A \lif !B$
  \end{prooftree}
  ફરી આગમન-પરિકલ્પના કહે છે કે આધાર-સિક્વન્ટ પ્રામાણ્ય છે; આપણે
  ફલિત-સિક્વન્ટ પણ પ્રામાણ્ય છે તેમ બતાવવા માગીએ છીએ. કોઈ મનસ્વી
  \iftag{FOL}{$\Struct{M}$}{$\pAssign{v}$} લો. સિક્વન્ટ
  $!A, \Gamma \Sequent \Delta, !B$ પ્રામાણ્ય હોવાથી નીચેનામાંથી ઓછામાં
  ઓછો એક કિસ્સો સધાય છે: (a)
  $\iftag{FOL}{\Sat/{M}}{\pSat/{v}}{!A}$, (b)
  $\iftag{FOL}{\Sat{M}}{\pSat{v}}{!B}$, (c) કોઈ $~!C \in \Gamma$ માટે
  $\iftag{FOL}{\Sat/{M}}{\pSat/{v}}{!C}$, અથવા (d) કોઈ
  $!C \in \Delta$ માટે $\iftag{FOL}{\Sat{M}}{\pSat{v}}{!C}$.
  કિસ્સા (a) અને (b)માં $\iftag{FOL}{\Sat{M}}{\pSat{v}}{!A \lif !B}$,
  અને તેથી એવું $!C \in \Delta, !A \lif !B$ છે કે
  $\iftag{FOL}{\Sat{M}}{\pSat{v}}{!C}$. કિસ્સા (c)માં કોઈ
  $!C \in \Gamma$ માટે $\iftag{FOL}{\Sat/{M}}{\pSat/{v}}{!C}$.
  કિસ્સા (d)માં કોઈ $!C \in \Delta$ માટે
  $\iftag{FOL}{\Sat{M}}{\pSat{v}}{!C}$. દરેક કિસ્સામાં
  \iftag{FOL}{$\Struct{M}$}{$\pAssign{v}$} એ
  $\Gamma \Sequent \Delta, !A \lif !B$ને સંતોષે છે. આ
  \iftag{FOL}{$\Struct{M}$}{$\pAssign{v}$} મનસ્વી હતું, તેથી
  $\Gamma \Sequent \Delta, !A \lif !B$ પ્રામાણ્ય છે. \iftag{FOL}{%
\item છેલ્લું અનુમાન \LeftR{\lforall} છે. ત્યારે એવું !!a{formula}~$!A(x)$
  અને બંધ પદ~$t$ છે કે $\pi$નો અંત આ પ્રમાણે છે:
  \begin{prooftree}
    \AxiomC{}
    \Deduce$!A(t), \Gamma \fCenter \Delta$
    \RightLabel{\LeftR{\lforall}}
    \UnaryInf$\lforall[x][!A(x)], \Gamma \fCenter \Delta$
  \end{prooftree}
  આપણે ફલિત-સિક્વન્ટ $\lforall[x][!A(x)], \Gamma \Sequent \Delta$
  પ્રામાણ્ય છે તેમ બતાવવા માગીએ છીએ. કોઈ !!a{structure}~$\Struct M$ લો.
  આધાર-સિક્વન્ટ $!A(t), \Gamma \Sequent \Delta$ પ્રામાણ્ય હોવાથી કાં તો
  (a) $\Sat/{M}{!A(t)}$, અથવા (b) કોઈ $!C \in \Gamma$ માટે
  $\Sat/{M}{!C}$, અથવા (c) કોઈ $!C \in \Delta$ માટે $\Sat{M}{!C}$.
  કિસ્સા (a)માં, \olref[syn][sem]{prop:quant-terms} મુજબ જો
  $\Sat{M}{\lforall[x][!A(x)]}$ હોય, તો $\Sat{M}{!A(t)}$. પરંતુ
  $\Sat/{M}{!A(t)}$, તેથી $\Sat/{M}{\lforall[x][!A(x)]}$. કિસ્સા (b)
  અને (c)માં $\Struct{M}$ એ
  $\lforall[x][!A(x)], \Gamma \Sequent \Delta$ને પણ સંતોષે છે.
  $\Struct M$ મનસ્વી હતું, તેથી
  $\lforall[x][!A(x)], \Gamma \Sequent \Delta$ પ્રામાણ્ય છે.
\item છેલ્લું અનુમાન \RightR{\lexists} છે: અભ્યાસપ્રશ્ન.
\item છેલ્લું અનુમાન \RightR{\lforall} છે. ત્યારે એવું !!a{formula}~$!A(x)$
  અને એવું !!a{constant}~$a$ છે કે $\pi$નો અંત આ પ્રમાણે છે:
  \begin{prooftree}
    \AxiomC{}
    \Deduce$\Gamma \fCenter \Delta, !A(a)$
    \RightLabel{\RightR{\lforall}}
    \UnaryInf$\Gamma \fCenter \Delta, \lforall[x][!A(x)]$
  \end{prooftree}
  અહીં આઇગનચલ-શરત સંતોષાય છે, એટલે કે $a$ એ $!A(x)$, $\Gamma$ કે
  $\Delta$માં આવતું નથી. આગમન-પરિકલ્પના મુજબ છેલ્લા અનુમાનનું
  આધાર-સિક્વન્ટ પ્રામાણ્ય છે. તેનું ફલિત-સિક્વન્ટ પણ પ્રામાણ્ય છે તે
  બતાવવું છે; એટલે કે કોઈપણ !!{structure}~$\Struct M$ માટે કાં તો
  (a) $\Sat{M}{\lforall[x][!A(x)]}$, અથવા (b) કોઈ $!C \in \Gamma$ માટે
  $\Sat/{M}{!C}$, અથવા (c) કોઈ $!C \in \Delta$ માટે $\Sat{M}{!C}$.

  માનો કે $\Struct{M}$ કોઈ મનસ્વી !!{structure} છે. (b) અથવા (c) સધાય
  તો કામ પૂરું થયું; તેથી માનો કે બંનેમાંથી એકેય સધાતું નથી: દરેક
  $!C \in \Gamma$ માટે $\Sat{M}{!C}$ અને દરેક $!C \in \Delta$ માટે
  $\Sat/{M}{!C}$. આપણે (a), એટલે કે
  $\Sat{M}{\lforall[x][!A(x)]}$, બતાવવાનું છે.
  \olref[syn][ass]{prop:sat-quant} મુજબ દરેક ચલ-નિયુક્તિ~$s$ માટે
  $\Sat{M}{!A(x)}[s]$ બતાવવું પૂરતું છે. તેથી કોઈ મનસ્વી ચલ-નિયુક્તિ
  $s$ લો. એવી સંરચના~$\Struct{M'}$ લો જે $\Struct{M}$ જેવી જ છે, સિવાય
  કે $\Assign{a}{M'} = s(x)$. \olref[syn][ext]{cor:extensionality-sent}
  મુજબ કોઈપણ $!C \in \Gamma$ માટે $\Sat{M'}{!C}$, કારણ કે $a$ એ
  $\Gamma$માં આવતું નથી; અને કોઈપણ $!C \in \Delta$ માટે
  $\Sat/{M'}{!C}$. પરંતુ આધાર-સિક્વન્ટ પ્રામાણ્ય છે, તેથી
  $\Sat{M'}{!A(a)}$. $!A(a)$ વાક્ય હોવાથી
  \olref[syn][ass]{prop:sentence-sat-true} મુજબ
  $\Sat{M'}{!A(a)}[s]$. હવે $\varAssign{s}{s}{x}$ અને
  $s(x) = \Value{a}{M'}[s]$, કારણ કે આપણે $\Struct{M'}$ને એ જ રીતે
  વ્યાખ્યાયિત કરી છે. તેથી \olref[syn][ext]{prop:ext-formulas} લાગુ પડે
  છે અને $\Sat{M'}{!A(x)}[s]$ મળે છે. $a$ એ~$!A(x)$માં આવતું નથી,
  તેથી \olref[syn][ext]{prop:extensionality} મુજબ
  $\Sat{M}{!A(x)}[s]$. $s$ મનસ્વી હતું, એટલે દરેક ચલ-નિયુક્તિ માટે
  $\Sat{M}{!A(x)}[s]$ સાબિત થયું.
\item છેલ્લું અનુમાન \LeftR{\lexists} છે: અભ્યાસપ્રશ્ન.
}{}
\end{enumerate}
હવે બે આધાર-સિક્વન્ટ ધરાવતાં સંભવિત અનુમાનો જોઈએ.
\begin{enumerate}
\item છેલ્લું અનુમાન કટ છે. ત્યારે $\pi$નો અંત આ પ્રમાણે છે:
  \begin{prooftree}
    \AxiomC{}
    \Deduce$\Gamma \fCenter \Delta, !A$
    \AxiomC{}
    \Deduce$!A, \Pi \fCenter \Lambda$
    \RightLabel{\Cut}
    \BinaryInf$\Gamma, \Pi \fCenter \Delta, \Lambda$
  \end{prooftree}
  કોઈ \iftag{FOL}{!!a{structure}~$\Struct{M}$}{!!a{valuation}~$\pAssign{v}$}
  લો. આગમન-પરિકલ્પના મુજબ બંને આધાર-સિક્વન્ટો પ્રામાણ્ય છે, તેથી
  \iftag{FOL}{$\Struct{M}$}{$\pAssign{v}$} બંનેને સંતોષે છે. બે કિસ્સા
  પાડીએ: (a) $\iftag{FOL}{\Sat/{M}}{\pSat/{v}}{!A}$ અને
  (b) $\iftag{FOL}{\Sat{M}}{\pSat{v}}{!A}$. કિસ્સા (a)માં ડાબા
  આધાર-સિક્વન્ટને સંતોષવા માટે
  \iftag{FOL}{$\Struct{M}$}{$\pAssign{v}$} એ
  $\Gamma \Sequent \Delta$ને સંતોષવું જ પડે. પરંતુ ત્યારે તે
  ફલિત-સિક્વન્ટને પણ સંતોષે છે. કિસ્સા (b)માં જમણા આધાર-સિક્વન્ટને
  સંતોષવા માટે \iftag{FOL}{$\Struct{M}$}{$\pAssign{v}$} એ
  $\Pi \Sequent \Lambda$ને સંતોષવું જ પડે. ફરીથી,
  \iftag{FOL}{$\Struct{M}$}{$\pAssign{v}$} ફલિત-સિક્વન્ટને સંતોષે છે.
  \sourcecorrection{OLSEQ-006}{મૂળની \texttt{soundness.tex}માં કટના
  બીજા કિસ્સામાં જમણું આધાર-સિક્વન્ટ $\Pi \Sequent \Lambda$ હોવું
  જોઈએ ત્યાં ગણ-બાદબાકી $\Pi \setminus \Lambda$ લખાઈ છે. અહીં
  સાબિતીમાં વપરાતું સાચું સિક્વન્ટ પુનઃસ્થાપિત કર્યું છે.}
\item છેલ્લું અનુમાન \RightR{\land} છે. ત્યારે $\pi$નો અંત આ પ્રમાણે છે:
  \begin{prooftree}
    \AxiomC{}
    \Deduce$\Gamma \fCenter \Delta, !A$
    \AxiomC{}
    \Deduce$\Gamma \fCenter \Delta, !B$
    \RightLabel{\RightR{\land}}
    \BinaryInf$\Gamma \fCenter \Delta, !A \land !B$
  \end{prooftree}
  કોઈ \iftag{FOL}{!!a{structure}~$\Struct M$}{!!a{valuation}~$\pAssign{v}$}
  લો. જો \iftag{FOL}{$\Struct{M}$}{$\pAssign{v}$} એ
  $\Gamma \Sequent \Delta$ને સંતોષે, તો કામ પૂરું થયું. તેથી માનો કે તે
  તેને સંતોષતું નથી. આગમન-પરિકલ્પના મુજબ
  $\Gamma \fCenter \Delta, !A$ પ્રામાણ્ય હોવાથી
  $\iftag{FOL}{\Sat{M}}{\pSat{v}}{!A}$. એ જ રીતે,
  $\Gamma \Sequent \Delta, !B$ પ્રામાણ્ય હોવાથી
  $\iftag{FOL}{\Sat{M}}{\pSat{v}}{!B}$. તેથી
  $\iftag{FOL}{\Sat{M}}{\pSat{v}}{!A \land !B}$.
\item છેલ્લું અનુમાન \LeftR{\lor} છે: અભ્યાસપ્રશ્ન.
\item છેલ્લું અનુમાન \LeftR{\lif} છે. ત્યારે $\pi$નો અંત આ પ્રમાણે છે:
  \begin{prooftree}
    \AxiomC{}
    \Deduce$\Gamma \fCenter \Delta, !A$
    \AxiomC{}
    \Deduce$!B, \Pi \fCenter \Lambda$
    \RightLabel{\LeftR{\lif}}
    \BinaryInf$!A \lif !B, \Gamma, \Pi \fCenter \Delta, \Lambda$
  \end{prooftree}
  ફરી કોઈ
  \iftag{FOL}{!!a{structure}~$\Struct{M}$}{!!a{valuation}~$\pAssign{v}$}
  લો અને માનો કે \iftag{FOL}{$\Struct{M}$}{$\pAssign{v}$} એ
  $\Gamma, \Pi \Sequent \Delta, \Lambda$ને સંતોષતું નથી. આપણે
  $\iftag{FOL}{\Sat/{M}}{\pSat/{v}}{!A \lif !B}$ બતાવવાનું છે. જો
  \iftag{FOL}{$\Struct{M}$}{$\pAssign{v}$} એ
  $\Gamma, \Pi \Sequent \Delta, \Lambda$ને સંતોષતું ન હોય, તો તે
  $\Gamma \Sequent \Delta$ કે $\Pi \Sequent \Lambda$, એકેયને
  સંતોષતું નથી. $\Gamma \Sequent \Delta, !A$ પ્રામાણ્ય હોવાથી
  $\iftag{FOL}{\Sat{M}}{\pSat{v}}{!A}$. $!B, \Pi \Sequent \Lambda$
  પ્રામાણ્ય હોવાથી $\iftag{FOL}{\Sat/{M}}{\pSat/{v}}{!B}$. તેથી
  $\iftag{FOL}{\Sat/{M}}{\pSat/{v}}{!A \lif !B}$, જે જ બતાવવાનું હતું.
\end{enumerate}
\end{proof}

\tagprob{FOL}
\begin{prob}
\olref[fol][seq][sou]{thm:sequent-soundness}ની સાબિતી પૂરી કરો.
\end{prob}
\tagendprob

\tagprob{notFOL}
\begin{prob}
\olref[pl][seq][sou]{thm:sequent-soundness}ની સાબિતી પૂરી કરો.
\end{prob}
\tagendprob

\begin{cor}
\ollabel{cor:weak-soundness}
જો $\Proves !A$ હોય, તો $!A$ \iftag{FOL}{પ્રામાણ્ય}{પુનરુક્તિ} છે.
\end{cor}

\begin{cor}
\ollabel{cor:entailment-soundness}
જો $\Gamma \Proves !A$ હોય, તો $\Gamma \Entails !A$.
\end{cor}

\begin{proof}
જો $\Gamma \Proves !A$ હોય, તો કોઈ સાન્ત ઉપગણ
$\Gamma_0 \subseteq \Gamma$ માટે $\Gamma_0 \Sequent !A$ની
!!a{derivation} છે. \olref{thm:sequent-soundness} મુજબ દરેક
\iftag{FOL}{!!{structure}~$\Struct{M}$}{!!{valuation}~$\pAssign{v}$}
કાં તો કોઈ $!B \in \Gamma_0$ને અસત્ય કરે છે અથવા $!A$ને સત્ય કરે છે.
તેથી જો $\iftag{FOL}{\Sat{M}}{\pSat{v}}{\Gamma}$ હોય, તો
$\iftag{FOL}{\Sat{M}}{\pSat{v}}{!A}$ પણ છે.
\end{proof}

\begin{cor}
\ollabel{cor:consistency-soundness}
જો $\Gamma$ સંતોષ્ય હોય, તો તે સુસંગત છે.
\end{cor}

\begin{proof}
પ્રતિપક્ષી વિધાન સાબિત કરીએ. માનો કે $\Gamma$ સુસંગત નથી. ત્યારે કોઈ
સાન્ત $\Gamma_0 \subseteq \Gamma$ અને
$\Gamma_0 \Sequent \quad$ની !!a{derivation} છે.
\olref{thm:sequent-soundness} મુજબ $\Gamma_0 \Sequent \quad$ પ્રામાણ્ય
છે. બીજા શબ્દોમાં, દરેક
\iftag{FOL}{!!{structure}~$\Struct{M}$}{!!{valuation}~$\pAssign{v}$}
માટે એવું $!C \in \Gamma_0$ છે કે
$\iftag{FOL}{\Sat/{M}}{\pSat/{v}}{!C}$; અને
$\Gamma_0 \subseteq \Gamma$ હોવાથી તે $!C$ એ~$\Gamma$માં પણ છે. તેથી
કોઈ \iftag{FOL}{$\Struct{M}$}{$\pAssign{v}$} એ~$\Gamma$ને સંતોષતું નથી,
અને $\Gamma$ સંતોષ્ય નથી.
\end{proof}

\end{document}
